% =============================================================================
% The Reproducibility Crisis in Artificial Intelligence:
% Statistical Biases, Benchmark Limitations, and Methodological Solutions
% =============================================================================
\documentclass[12pt, a4paper]{article}

% --- Encoding and language ---------------------------------------------------
\usepackage[utf8]{inputenc}
\usepackage[T1]{fontenc}
\usepackage[english]{babel}

% --- Typography and layout ---------------------------------------------------
\usepackage{lmodern}
\usepackage{microtype}
\usepackage[margin=2.5cm]{geometry}
\usepackage{setspace}
\onehalfspacing

% --- Bibliography ------------------------------------------------------------
\usepackage[numbers,sort&compress]{natbib}
\bibliographystyle{unsrtnat}

% --- Cross-references and hyperlinks -----------------------------------------
\usepackage[
  colorlinks  = true,
  linkcolor   = blue!70!black,
  citecolor   = blue!70!black,
  urlcolor    = blue!60!black,
  bookmarks   = true,
  pdfauthor   = {},
  pdftitle    = {The Reproducibility Crisis in Artificial Intelligence},
  pdfsubject  = {Reproducibility, Machine Learning, Benchmarking},
]{hyperref}
\usepackage[nameinlink]{cleveref}

% --- Miscellaneous -----------------------------------------------------------
\usepackage{booktabs}
\usepackage{enumitem}
\usepackage{xcolor}

% =============================================================================
\title{%
  The Reproducibility Crisis in Artificial Intelligence:\\[4pt]
  \large Statistical Biases, Benchmark Limitations,\\
  and Methodological Solutions%
}

\author{}
\date{\today}

% =============================================================================
\begin{document}
\maketitle
\thispagestyle{empty}

% -----------------------------------------------------------------------------
% Abstract
% -----------------------------------------------------------------------------
\begin{abstract}
\noindent
Empirical research in Artificial Intelligence (AI) depends heavily on
data-driven experiments, benchmarks, and complex software--hardware
ecosystems.  This dependence creates distinctive risks for
reproducibility and reliable inference that fall into three broad
categories: \emph{(i)}~statistical biases---including p-hacking
analogues that arise from the many degrees of freedom in model and
training decisions, as well as selection and publication biases;
\emph{(ii)}~systematic over-adaptation to benchmarks, encompassing
adaptive misuse of holdout sets and Goodhart effects; and
\emph{(iii)}~methodological error classes such as overfitting, dataset
leakage, duplicate samples, and label errors, all of which can
substantially distort reported performance figures.

Meta-research and community programmes have documented persistent
shortcomings in the reporting of code, data, and experimental
configurations for years, yet they also demonstrate that institutional
interventions---checklists, artifact reviews, reproducibility
challenges---can trigger measurable behavioural change
\cite{Pineau2021,NeurIPSChecklist2023}.

This paper synthesises the scientific literature on the
``reproducibility crisis'' in AI along statistical, infrastructural, and
governance-related dimensions \cite{Goodman2016,NationalAcademies2019}.
Furthermore, it presents a solution framework that combines Open Science
practices, technical infrastructure (versioning, provenance, packaging),
dataset governance (documentation standards), and automated citation and
evidence verification to raise the validity of empirical AI results and
to reduce the perverse incentives of the benchmark ecosystem.

\medskip\noindent
\textbf{Keywords:}
reproducibility, replication, p-hacking, publication bias, dataset
leakage, benchmarking, adaptive data analysis, Open Science, data and
software citation, provenance, versioning, evidence and citation
verification.
\end{abstract}

\newpage
\tableofcontents
\newpage

% =============================================================================
\section{Introduction}
\label{sec:introduction}
% =============================================================================

Reproducibility and replication are cornerstones of cumulative science,
yet in practice they are compromised by statistical, social, and
infrastructural factors across many disciplines.  The very terminology
remains incompletely standardised; a key finding of meta-science is that
precise conceptual work---for instance, distinguishing
\emph{methods reproducibility}, \emph{results reproducibility}, and
\emph{inferential reproducibility}---is a prerequisite for
operationalising crisis diagnoses and evaluating interventions at all
\cite{Goodman2016}.

In AI---including machine learning (ML), deep learning, and
data-intensive subfields such as natural language processing (NLP),
computer vision, and reinforcement learning (RL)---the classical
reproducibility risks are amplified by three structural properties of
the field:

\begin{enumerate}[leftmargin=*]
  \item \textbf{Stochastic training and fragile baselines.}
    Central empirical claims frequently rest on small performance
    differences on benchmarks, yet these differences are subject to
    substantial fluctuation from stochastic training processes,
    non-deterministic execution environments, and hyperparameter search
    \cite{Henderson2018,Bouthillier2021}.

  \item \textbf{Complex, under-documented artefacts.}
    Research artefacts---code, data, models, preprocessing pipelines,
    dependencies, hardware and driver versions---are highly complex and
    often insufficiently documented or archived, impeding both
    repetition and independent verification
    \cite{Peng2011,Stodden2016,Gundersen2020}.

  \item \textbf{Benchmark-driven incentives and Goodhart dynamics.}
    Benchmarks and leaderboards create powerful incentive structures.
    When a metric becomes a target, it can lose its validity as a
    measurement instrument (Goodhart effect), particularly through
    adaptive overuse of static test sets and through metrics that lend
    themselves to gaming \cite{Hsia2024,Raji2021}.
\end{enumerate}

The scientific community has responded over the past several years with
programmatic and infrastructural measures: conferences have implemented
checklists and code policies; artifact review processes and
reproducibility challenges create new forms of publication and
recognition; Open Science and FAIR principles target the findability,
reusability, and machine-actionability of research objects
\cite{NeurIPSChecklist2023,Wilkinson2016,Nosek2015}.  At the same time,
the literature shows that technical solutions (versioning, provenance,
packaging, archivability) have limited effectiveness without adequate
governance and without statistical methodological discipline, because
biases can already arise in the research logic itself---in hypothesis
generation, model selection, and reporting and publication dynamics
\cite{Kapoor2023,Munafo2017}.

The goal of this paper is a deep, source-based synthesis of the
reproducibility problem in AI along the following dimensions: statistical
foundations (p-hacking, publication bias), overfitting and dataset
leakage, benchmark optimisation problems, meta-research evidence,
Open Science initiatives, technical solutions (versioning, dataset
governance, automated citation verification), case studies, and a
concluding discussion that integrates these threads into a unified
solution framework.

% =============================================================================
\section{Statistical and Methodological Causes}
\label{sec:statistical}
% =============================================================================

Before examining AI-specific pathologies, it is useful to establish the
broader statistical context in which the reproducibility crisis is
rooted.  Two mutually reinforcing mechanisms---flexible analysis
decisions and selective publication---generate systematic distortions
that inflate the apparent strength and reliability of empirical findings.
This section traces these mechanisms from their general-science origins
into their AI-specific manifestations, and then turns to the
methodological error classes (overfitting, leakage, contamination) that
amplify the problem.

% -----------------------------------------------------------------------------
\subsection{Statistical Foundations: p-Hacking Analogues, Multiple
            Degrees of Freedom, and Publication Bias}
\label{sec:phacking}
% -----------------------------------------------------------------------------

The classical p-hacking debate demonstrates that flexible analysis
decisions (``researcher degrees of freedom'') and selective reporting can
render significant results more probable even in the absence of robust
effects.  This phenomenon is well documented across many empirical
sciences and is discussed as a central mechanism behind false-positive
findings \cite{Simmons2011,Ioannidis2005,Head2015}.

In AI, a functional analogue manifests not as manipulation of individual
p-values but as optimisation across many degrees of freedom: choice of
architecture, data augmentation, preprocessing, optimiser, training
budget, random seed, hyperparameter space, and selection criterion
(e.g., best single run versus mean across seeds).  When only the best
configuration is subsequently reported, a selection bias emerges that is
formally equivalent to the multiple-testing problem
\cite{Bouthillier2021,Henderson2018}.

Several studies emphasise that pure ``determinism'' reproducibility
(same code, same seed, identical environment) is insufficient to secure
inferential robustness.  Instead, variance from multiple sources---data
sampling, initialisation, hyperparameter choice, implementation
details---must be modelled and integrated into the interpretation of
results \cite{Bouthillier2021,Hagmann2023}.  This demand is closely
related to statistical guidelines for correct algorithm-comparison
inference, such as appropriate tests over multiple datasets or
experimental conditions \cite{Demsar2006,Dror2018}.

Publication bias (the ``file drawer problem'') further compounds the
distortion: when ``positive'' or ``state-of-the-art'' results enjoy a
higher publication probability than null or negative results, effect
sizes and progress rates are systematically overestimated.  This is well
established in meta-analyses and population-based studies of publication
dynamics \cite{Franco2014,Fanelli2010}.  For AI, the analogy is
discussed primarily as an incentive problem: benchmarks, competitions,
and community norms frequently reward incremental performance gains,
whereas replication studies, negative results, or ``ablation null
findings'' receive lesser recognition.  Open Science manifestos and
recommendations for reforming incentive and publication systems address
these dynamics explicitly \cite{Munafo2017,Nosek2015}.

Methodological counter-reactions from the philosophy of science and
meta-science literature include, among others, preregistration,
Registered Reports (review before data analysis or before knowledge of
results), and stronger transparency requirements.  Registered Reports
have been established as a publication format to reduce
results-orientation and to tie the publication decision more closely to
methodology than to significance \cite{Chambers2013}.  In AI, direct
transfer is not trivial (e.g., iterative development, high experimental
costs), yet the core principle---separating the analysis plan from the
results---remains conceptually highly relevant as a countermeasure
against selective reporting \cite{Goodman2016}.

% -----------------------------------------------------------------------------
\subsection{Overfitting and Dataset Leakage: Forms, Causes, and
            Empirical Evidence}
\label{sec:leakage}
% -----------------------------------------------------------------------------

The general statistical biases described above are compounded by a set
of methodological error classes specific to the data-centric nature of
ML research.  Overfitting in AI is not merely a theoretical concept of
the learning curve; it is tightly coupled with benchmark culture,
data-pipeline design, and evaluation practices.

A central error class is \textbf{dataset leakage}: information from test
or validation data enters training or tuning processes, whether
implicitly or explicitly---for instance through incorrect splitting,
preprocessing on the full dataset, leakage via temporal indices, or
meta-feature engineering \cite{Kapoor2023}.  A comprehensive,
cross-domain study argues that leakage in ML-based science is
widespread and has in numerous cases led to ``wildly overoptimistic
conclusions''; it further proposes standardised ``Model Info Sheets'' to
systematically audit leakage risks \cite{Kapoor2023}.

In domains with highly structured data (e.g., medical imaging), even the
\emph{granularity} of the split---patient/study/volume level versus
image/slice level---can determine whether seemingly outstanding
performance is real or artificially inflated by leakage.  A detailed
analysis of optical coherence tomography (OCT) image classification
shows that inappropriate splits (with overlap between training and test
sets) can markedly inflate accuracy figures, reporting large performance
differences between correct and incorrect splitting strategies
\cite{Tampu2022}.

Alongside ``classical'' leakage, a second, structurally related problem
class exists: \textbf{duplicates and near-duplicates} between training
and test sets.  For heavily used image datasets, it has been shown that
a non-trivial proportion of test examples have duplicates in the
training set; this can distort evaluation comparability because models
effectively (partially) train on test data or benefit from memorisation.
A study on CIFAR-10/100 identifies near-duplicates in the test set,
replaces them with fresh examples (``ciFAIR''), and observes significant
accuracy drops---a direct indication that benchmark values can be
systematically influenced by data overlap \cite{Barz2020}.

A third problem source consists of \textbf{label errors in test sets}.
A systematic characterisation of label errors across several popular
benchmarks not only identifies the presence of errors but also
quantifies their magnitude (e.g., error proportions in commonly used
test/validation splits) and discusses the consequences for benchmark
stability and model selection \cite{Northcutt2021}.  Even when rankings
may in some cases remain robust, the normative problem persists: if the
measurement basis itself is flawed, absolute performance values and fine
distinctions---which frequently drive publication decisions---are less
informative \cite{Northcutt2021}.

Finally, the proliferation of very large pretrained models (particularly
large language models, LLMs) has introduced a new form of leakage:
\textbf{benchmark contamination through pretraining data}.  When
benchmark items or semantically close paraphrases are contained in
web-scale corpora, models can memorise portions of the evaluation data.
Several studies develop detection methods, taxonomise contamination
types, and propose documentation artefacts (e.g., ``Benchmark
Transparency Cards'') to improve training-data transparency and
comparability \cite{XuR2024,XuC2024,Choi2025}.

% =============================================================================
\section{Limitations of Benchmarks and Optimisation Dynamics}
\label{sec:benchmarks}
% =============================================================================

The methodological problems identified in the previous section do not
exist in isolation; they interact with the \emph{institutional}
structure of how AI research is evaluated.  Benchmarks simultaneously
serve as measurement instruments and as steering mechanisms for research
attention and resource allocation.  This dual role creates optimisation
dynamics that can undermine the very validity benchmarks are designed to
ensure.

% -----------------------------------------------------------------------------
\subsection{Benchmark Optimisation as an Adaptive Inference Problem}
\label{sec:adaptive}
% -----------------------------------------------------------------------------

In this respect, benchmarks resemble adaptive data analysis: hypotheses,
models, and pipeline decisions are iteratively generated based on
benchmark feedback.  The statistics literature shows that repeated,
adaptive use of a holdout dataset can undermine the validity of
classical inference assumptions, because the ``holdout'' becomes
implicitly integrated into the decision space \cite{Dwork2015}.

In ML competitions and leaderboards this problem is operationally
particularly visible: participants can overfit their solutions to a
public leaderboard through repeated submission.  Theoretical work models
this risk and proposes mechanisms designed to reduce leaderboard
overfitting by disclosing holdout information only in a restricted,
adaptation-resilient manner \cite{Blum2015}.  This perspective transfers
to scientific benchmarks, where iterative tuning and ``leaderboard
chasing'' are typical research practices, especially when improvements
on the order of a few tenths of a percentage point suffice for
publication \cite{Bouthillier2021}.

% -----------------------------------------------------------------------------
\subsection{Goodhart Effects, ``General'' Benchmarks, and External
            Validity}
\label{sec:goodhart}
% -----------------------------------------------------------------------------

When a metric becomes a target, it can lose its informational value
(Goodhart effect).  In AI, this manifests in several forms:
\emph{(i)}~models learn spurious correlations and shortcuts instead of
intended capabilities; \emph{(ii)}~task metrics can be systematically
``gamed'' without improving semantic quality; \emph{(iii)}~benchmark
design can unintentionally steer development toward metrically
favourable but scientifically or socially less valid properties.
Specific analyses demonstrate Goodhart mechanisms even within
specialised benchmark families (e.g., explanation benchmarks), where
scores can be substantially increased without genuine improvement in
explanatory substance \cite{Hsia2024}.

The critique of ``general'' benchmarks goes further, arguing that a
small canon of prominent datasets in various subfields can be
misinterpreted as a proxy for ``general progress.''  A positional
analysis describes such general benchmarks as structurally incompatible
with strong claims of general capability, because benchmarks are
necessarily under-specified and the reinterpretation of state-of-the-art
(SOTA) performance on a benchmark as ``general intelligence'' is
epistemically unjustified \cite{Raji2021}.

The external-validity conflict is further exacerbated by
\textbf{dataset shift}: when the data distribution differs between
training and test sets---or between benchmark and deployment---benchmark
performance can only limitedly predict real-world performance.  The
problem has been discussed for years under the headings ``dataset
shift'' and ``covariate shift,'' including consequences for evaluation
design and robust modelling \cite{QuinoneroCandela2009}.  Empirical
benchmark research replicates this insight with new test sets or ``fresh
samples'': when models are adapted to a static test-set canon over
years, performance drops on newly drawn, but ``equivalently intended,''
test sets can be interpreted as indicators of benchmark overfitting
\cite{Recht2019}.

% -----------------------------------------------------------------------------
\subsection{Benchmark Design Beyond Leaderboards: Suites Instead of
            Single-Number Rankings}
\label{sec:suites}
% -----------------------------------------------------------------------------

A growing strand of the literature argues for transitioning from
leaderboards to \textbf{benchmark suites}: instead of optimising ``one
number'' (and hence one ranking), different metrics, usage scenarios,
and trade-offs should be made explicit.  This is articulated
particularly clearly in fairness evaluation: ``fairness'' is
conceptually plural and application-dependent; a leaderboard that
reduces fairness to a single score produces target distortions and can
obscure potential harms.  Benchmark suites, by contrast, are intended to
make different perspectives and measurement dimensions visible, thereby
increasing both scientific validity and governance-related
decision-making capacity \cite{Wang2024}.

% =============================================================================
\section{Evidence from Meta-Research and Open Science Initiatives}
\label{sec:meta}
% =============================================================================

Having established the theoretical mechanisms through which biases and
benchmark dynamics compromise AI research, we now turn to the empirical
evidence.  Meta-research quantifies the extent of the problem, while
Open Science initiatives represent the community's institutional
response.

% -----------------------------------------------------------------------------
\subsection{Meta-Research: Findings on Documentation and Reproduction
            Practice in AI}
\label{sec:metaresearch}
% -----------------------------------------------------------------------------

Meta-research in AI consistently shows that pure results reports
frequently do not contain sufficient information to reliably reproduce
experiments.  A documentation-oriented analysis spanning several major AI
conferences (sample across multiple years) indicates, for example, that
only a small proportion of the examined papers share centrally relevant
artefacts (including code, results/outputs, hypotheses/predictions) and
that experimental details (dependencies, setup parameters) can be
severely under-documented \cite{Gundersen2020}.  These findings should
not be read as mere ``compliance'' criticism; they represent an
\emph{epistemic} risk: when essential pipeline decisions are not auditable,
it is difficult to distinguish whether a result rests on a robust
methodological innovation or on unreported, potentially data-specific or
chance-driven factors \cite{Goodman2016}.

Meta-research on variance and benchmarking additionally reveals that
benchmark values are often reported as point estimates despite being
subject to stochastic and pipeline-related fluctuations.  Work on
``variance in machine learning benchmarks'' models the benchmark
pipeline and argues that common comparison practices underestimate
variance, thereby producing a distorted picture of methodological
superiority \cite{Bouthillier2021}.  For subfields such as reinforcement
learning, this problem is formulated particularly sharply:
non-determinism in environments and intrinsic variance can render
purported SOTA improvements difficult to interpret without significance
metrics and robust reporting standards \cite{Henderson2018}.

Finally, meta-research shows that not only \emph{results reproduction}
but also \emph{inferential reproducibility} (the validity of
conclusions across groups of tasks/datasets) constitutes a distinct
target criterion: rather than ``determinising variance away'' (e.g., by
fixing a single seed), the proposal is to deliberately vary variance
sources and to model them statistically (e.g., mixed-effects models) in
order to obtain generalisable statements \cite{Hagmann2023}.

% -----------------------------------------------------------------------------
\subsection{Open Science Initiatives: Checklists, Artifact Reviews, and
            Community Programmes}
\label{sec:openscience}
% -----------------------------------------------------------------------------

At the institutional level, several initiatives have emerged in recent
years to raise reproducibility in AI.  A central example is the
introduction of \textbf{paper checklists} by major ML conferences, which
query best practices regarding reproducibility, transparency, ethics,
and societal impact; the guidelines emphasise that the checklist must not
be removed and functions as a component of the submission
\cite{NeurIPSChecklist2023}.

An empirically well-documented intervention case is a conference
\textbf{reproducibility programme} that combines a code-submission
policy, a reproducibility challenge, and a checklist.  The programme
report describes, among other things, a marked increase in voluntary
code submissions (from below 50\% to nearly 75\% within one year) and
further reports that a portion of reviewers rated the checklist as
helpful.  At the same time, it emphasises that robust evidence on
long-term quality effects remains limited, underscoring the need for
further meta-evaluations of such interventions \cite{Pineau2021}.

Beyond conference mechanisms, \textbf{artifact review and badging
programmes} (particularly in the computing ecosystem) define criteria
under which artefacts are evaluated as repeatable/reproducible and how
this is made visible.  Such frameworks aim to close the ``artefact gap''
between paper and executable research and to create incentives for clean
packaging and documentation practices \cite{ACMArtifactReview2020}.

Open Science is further supported by broad, cross-disciplinary
principles and manifestos that formulate transparency, openness, and
reproducibility as systemic goals.  Key reference points include
\emph{(i)}~recommendations for promoting an open research culture
\cite{Nosek2015}, \emph{(ii)}~guidelines for the reproducibility of
computational methods, particularly code and data provision
\cite{Stodden2016,Peng2011}, and \emph{(iii)}~FAIR principles, which
emphasise the machine-actionability of research objects
\cite{Wilkinson2016}.

% =============================================================================
\section{Technical and Organisational Solutions}
\label{sec:solutions}
% =============================================================================

The preceding sections have established that the reproducibility crisis
in AI arises from an interplay of statistical, benchmark-related, and
documentation deficits.  Addressing these deficits requires not only
cultural and procedural change but also concrete technical
infrastructure and governance mechanisms.  This section reviews solutions
along three axes: versioning and executable reproduction, dataset
governance, and automated verification.

% -----------------------------------------------------------------------------
\subsection{Versioning, Provenance, and Executable Reproduction}
\label{sec:versioning}
% -----------------------------------------------------------------------------

Technical reproducibility in AI requires the controlled management of
code, data, model artefacts, and execution environments.  Data and code
versioning systems aim to make states (snapshots) of data and pipelines,
as well as their dependencies, traceable.

One example is \textbf{DataLad}, a distributed system for the joint
management of code and data that connects Git-based workflows with
large-scale data logistics and is explicitly oriented toward
reproducible research and FAIR compatibility \cite{Halchenko2021}.  For
Data Version Control (DVC), empirical analyses of usage in open-source
ML development characterise DVC as a tool for versioning and
reproducibility of ML artefacts \cite{Pacheco2024}.

A second core building block is the \textbf{long-term archiving and
referenceability of software sources}.  Initiatives for comprehensive
source-code archiving argue that persistent identifiers and systematic
archiving can address central reproducibility barriers by reducing
``link rot'' and the problem of untraceable code states
\cite{DiCosmo2018}.

Third, \textbf{packaging tools and containerised reproduction
environments} become relevant: reproducibility frequently fails due to
``dependency hell.''  Tools for packaging experiments extract
dependencies and enable cross-platform re-execution; corresponding
systems have been positioned at data-management and database conferences
as reproducibility infrastructure \cite{Chirigati2016,Rampin2018}.
Complementary platform approaches bundle code, data, dependencies, and
metadata in a ``capsule,'' check executability, and assign DOIs for
citability and archiving; this is described in the literature as a way
to reduce review and reproduction burden \cite{Cheifet2021}.

% -----------------------------------------------------------------------------
\subsection{Dataset Governance: Documentation, Quality Assurance, and
            Citation Practices}
\label{sec:governance}
% -----------------------------------------------------------------------------

Many reproducibility problems are data-driven.  Consequently, a research
line has emerged that understands dataset and model documentation as a
governance instrument.

\textbf{Datasheets for Datasets} propose standardised documents that
make the motivation, composition, collection process, recommended use,
and limitations of datasets transparent; in peer-reviewed form, such
standards have been prominently published \cite{Gebru2021}.  For NLP,
\textbf{Data Statements} exist as a professional practice format that
systematically describes data provenance, sampling, annotation, and
potential bias consequences to enable better science and bias mitigation
\cite{Bender2018}.  For models, comparable documentation formats
(\textbf{Model Cards}) have been proposed that disclose performance
profiles, deployment contexts, and limitations \cite{Mitchell2019}.

A leakage-specific extension of this governance idea takes the form of
\textbf{Model Info Sheets} as a checklist format that addresses concrete
leakage types and thereby supports the detection of methodological
errors prior to publication.  Empirical evidence from ML-based science
suggests that such systematic audit frameworks are necessary because
leakage errors are frequently not discernible from paper text alone
\cite{Kapoor2023}.

\textbf{Citation practices} constitute a further governance element
because they govern reusability, attribution, and machine-linkability of
research objects.  For data, a Joint Declaration has been formulated
that describes data citation as part of reproducible scholarship and
calls for machine- and human-readable citations \cite{FORCE11_2014}.
For software, software citation principles have been published in
peer-reviewed form, defining specificity (version), persistence, and
attribution as core elements \cite{Smith2016}.  In practice, these goals
are supported by formats such as the Citation File Format (CFF), in
which citable metadata are stored in a machine-readable manner
\cite{Druskat2021}.

Finally, \textbf{research packaging standards} (e.g., RO-Crate) address
the problem that paper text is often incomplete or ambiguous for
automated systems.  RO-Crate is described as a lightweight,
community-driven packaging format that bundles research artefacts and
metadata in a machine-readable fashion, thereby supporting FAIR goals
and reproducibility \cite{SoilandReyes2022}.

% -----------------------------------------------------------------------------
\subsection{Automated Citation and Evidence Verification Systems}
\label{sec:verification}
% -----------------------------------------------------------------------------

With the growth of scientific literature and the advent of AI-generated
content, the question of whether citations actually support the
assertions they accompany becomes increasingly pressing.  The connection
to the preceding infrastructure is direct: even the best versioning and
documentation frameworks provide limited value if the textual claims in
papers do not accurately reference the artefacts they are supposed to
describe.

The NLP literature has formalised benchmarks and tasks for this purpose:
\textbf{scientific claim verification} establishes datasets in which
claims are linked to evidence texts and annotated as
``supports/refutes,'' enabling the evaluation of systems for evidence
retrieval and reasoning \cite{Wadden2020}.  Complementary subtasks such
as \textbf{cite-worthiness detection} (where is a citation needed?) and
\textbf{citation intent classification} have been developed, forming
structural precursors to automated citation quality assurance
\cite{Wright2021}.

More recent benchmarks explicitly evaluate whether language models can
correctly attribute scientific claims and discuss the question of
whether systems can reliably perform the mapping ``claim $\to$ source.''
Such work positions automated attribution both as a means against
hallucinations and as a building block for future verification systems
\cite{Press2024}.  At the system level, proposals exist for
\textbf{citation verification} through full-text analysis with
classification taxonomies (e.g., supported/unsupported) and provision of
evidence snippets; corresponding systems argue that scalable citation
checking is integrity-relevant, especially for peer review and
AI-generated content \cite{SemanticCite2025}.

For the reproducibility crisis in AI, this line of work is particularly
relevant for two reasons.  First, automated citation checking can
\emph{enforce} the referencing of concrete artefacts (dataset versions,
code snapshots, model hashes) by detecting missing or incorrect
references and suggesting corrections.  Second, it can serve as a
quality-assurance mechanism in review pipelines to reduce the
discrepancy between textual assertions and actual
evidence/artefacts.  Both roles are explicitly formulated as target
directions in the literature---integrity, peer-review relief, and
evidence-based justification \cite{SemanticCite2025}.

% =============================================================================
\section{Case Studies}
\label{sec:casestudies}
% =============================================================================

The mechanisms and solutions discussed so far can be illustrated through
concrete cases drawn from different AI subfields.  Each case study
highlights a different facet of the reproducibility crisis and connects
back to the theoretical framework established in the preceding sections.

% -----------------------------------------------------------------------------
\subsection{Reinforcement Learning: Variance, Non-Determinism, and
            Inference Uncertainty}
\label{sec:case_rl}
% -----------------------------------------------------------------------------

An influential analysis of deep reinforcement learning problematises the
fact that non-deterministic environments and methodological variance can
render reported performance gains difficult to interpret when
significance measures, multiple runs, and standardised reporting
practices are lacking.  The finding is not merely an engineering
problem; it is a \emph{statistical inference problem}: without variance
modelling, it remains unclear whether differences are robust or
chance-driven \cite{Henderson2018}.  This case exemplifies the
p-hacking analogue discussed in \cref{sec:phacking}: researchers may
select the best seed or environment configuration, generating an
illusory impression of progress.

% -----------------------------------------------------------------------------
\subsection{Computer Vision Benchmarks: Near-Duplicates and Measurement
            Foundations}
\label{sec:case_cv}
% -----------------------------------------------------------------------------

The ciFAIR study demonstrates that near-duplicates between training and
test sets exist in heavily used benchmarks and that a ``purged'' test set
can reveal considerable performance losses.  The study interprets this
as an indication that benchmark progress can be partly facilitated by
memorisation and that benchmark design and dataset hygiene are central
components of scientific validity \cite{Barz2020}.  Complementary work
on label errors in test sets shows that even test data cannot be
presupposed as ``ground truth'' and that label quality can influence the
stability of benchmark comparisons \cite{Northcutt2021}.  Together,
these findings reinforce the leakage and data-quality concerns raised in
\cref{sec:leakage}.

% -----------------------------------------------------------------------------
\subsection{LLM Evaluation: Contamination and Dynamic Benchmarks}
\label{sec:case_llm}
% -----------------------------------------------------------------------------

Several studies document that benchmark contamination---overlap between
evaluation data and pretraining corpora---can inflate performance
claims and undermine fair comparisons; proposed remedies include
detection pipelines, survey systematisations, and transparency artefacts
\cite{XuR2024,XuC2024}.  At the same time, more recent analyses show
that the effect of contamination on generative benchmarks is causally
modulated by parameters such as sampling temperature and
solution/response length; this means that ``contamination'' must be
modelled not as a binary problem but as a context-dependent source of
bias \cite{Schaeffer2025}.

In response, \textbf{contamination-limited benchmarks} have emerged
that use frequently updated, objectively scorable tasks to reduce
contamination risk and to increase the reliability of LLM evaluation
\cite{LiveBench2024}.  This development represents a practical
implementation of the benchmark-suite philosophy discussed in
\cref{sec:suites}.

% -----------------------------------------------------------------------------
\subsection{ML-Based Science: Leakage as a Reproduction Error}
\label{sec:case_mlscience}
% -----------------------------------------------------------------------------

The leakage meta-study in ML-based science demonstrates that leakage is
not merely an internal ML error class but can have
reproducibility-critical consequences for entire application domains.  In
a reproduction study on civil-war prediction, it shows that the
purported superiority of complex ML models disappears upon correction
of leakage, thereby invalidating central conclusions of the original
literature \cite{Kapoor2023}.  This case underscores the broader point
made throughout this paper: performance claims that rest on compromised
evaluation pipelines can propagate unchecked through the scientific
record if reproducibility infrastructure and governance are absent.

% =============================================================================
\section{Discussion}
\label{sec:discussion}
% =============================================================================

The evidence assembled in this paper paints a consistent picture: the
``reproducibility crisis'' in AI is neither exclusively a problem of
``missing code links'' nor solely a problem of ``poor scientific
culture.''  Rather, it emerges from the interplay of statistical
selection logic, benchmark incentive mechanisms, and the infrastructural
complexity of computational experiments \cite{Goodman2016,Kapoor2023}.

A central insight is that AI-typical degrees of freedom can assume the
role of multiple hypothesis tests: configurations are iteratively
optimised, and the final report often condenses this process into a
single score.  The methodological consequence is that scientific
standards must be oriented more strongly toward measurement
uncertainties, variance components, and robust inference (e.g., multiple
seeds, appropriate test protocols, effect sizes instead of ``best
score'') \cite{Bouthillier2021,Hagmann2023}.

Benchmarks fulfil an economic function (standardisation,
comparability), but when misused they can generate epistemic damage:
adaptive overfitting, Goodhart effects, conflation of proxy targets with
``capabilities,'' and erosion of external validity through dataset
shift.  This does not mean that benchmarks are obsolete; it means that
their \emph{governance}---suites instead of leaderboards, regular
updating/rotation, leakage mitigation, quality control of labels and
duplicates---becomes a core component of scientific methodology
\cite{Wang2024}.

The intervention literature suggests that community mechanisms can be
effective in shifting norms (e.g., rising code submission rates,
establishment of checklists), but the evidence on long-term effects on
research quality remains incomplete.  This implies a \emph{meta-evaluation
imperative}: reproducibility programmes should themselves be understood
as evaluable interventions whose effectiveness (quality, robustness,
generalisability) is empirically tracked \cite{Pineau2021}.

Technically, a design principle emerges: reproducibility requires a
chain of \emph{referenceability} (persistent IDs), \emph{versioning}
(code/data/models), \emph{packaging} (executable environments),
\emph{documentation} (Datasheets / Data Statements / Model Cards / Info
Sheets), and finally \emph{auditability} (including automated evidence
and citation verification).  Each link addresses a different error
class: versioning reduces ``invisible'' drift; packaging reduces
execution barriers; documentation reduces interpretation and governance
deficits; automated citation checking reduces ``paper-to-artefact''
inconsistencies and raises the probability that claims actually rest on
traceable evidence \cite{Halchenko2021,SoilandReyes2022}.

% =============================================================================
\section{Conclusion}
\label{sec:conclusion}
% =============================================================================

The reproducibility crisis in AI is best understood as a
\emph{multi-layered validity problem}:
\begin{enumerate}[label=(\roman*), leftmargin=*]
  \item statistical biases through selection and reporting mechanisms;
  \item methodological error classes such as leakage, overfitting,
        duplicates, and erroneous labels;
  \item a benchmark economy with adaptive overuse and Goodhart effects;
        and
  \item infrastructural fragility through complex software and data
        dependencies \cite{Kapoor2023}.
\end{enumerate}

An effective solution approach must therefore simultaneously:
\begin{enumerate}[label=(\alph*), leftmargin=*]
  \item \textbf{modernise statistical inference practices}---variance
        modelling, robust tests, transparent reporting standards;
  \item \textbf{reform benchmark governance}---suites, updating, leakage
        control;
  \item \textbf{strengthen Open Science institutions}---checklists,
        artifact review, replication publication; and
  \item \textbf{expand technical infrastructure} toward versioned,
        packaged, and citable research objects---complemented by
        automated evidence and citation verification as a scaling
        mechanism for quality assurance \cite{Pineau2021}.
\end{enumerate}

None of these measures is sufficient in isolation.  Statistical reform
without infrastructure leaves researchers unable to share artefacts even
when willing; infrastructure without governance risks creating
technically reproducible but inferentially flawed research; governance
without statistical discipline cannot prevent the selection biases that
arise in research logic itself.  Only their \emph{combination}---a
coherent framework spanning statistics, benchmarks, institutions, and
technology---can credibly address the multi-dimensional nature of the
crisis and rebuild confidence in the empirical foundations of artificial
intelligence.

% =============================================================================
% Bibliography
% =============================================================================
\newpage
\bibliography{reproducibility_crisis_ai}

\end{document}
